
\chapter{2016年全国II卷（理）}

\section{单选题}

\begin{problem}
已知 \(z=(m+3)+(m-1)\i\) 在复平面内对应的点在第四象限，则实数 \(m\) 的取值范围是（\quad）

\choices
    {\((-3,1)\)}
    {\((-1,3)\)}
    {\((1,+\infty)\)}
    {\((-\infty,-3)\)}
\end{problem}

\begin{answer}
A
\end{answer}

\begin{problem}
已知集合 \(A=\{1,2,3\}\)，\(B=\{x\mid(x+1)(x-2)<0,\ x\in\Z\}\)，则 \(A\cup B\) 等于（\quad）

\choices
    {\(\{1\}\)}
    {\(\{1,2\}\)}
    {\(\{0,1,2,3\}\)}
    {\(\{-1,0,1,2,3\}\)}
\end{problem}

\begin{answer}
C
\end{answer}

\begin{problem}
已知向量 \(\bs a=(1,m)\)，\(\bs b=(3,-2)\)，且 \((\bs a+\bs b)\perp\bs b\)，则 \(m=\)（\quad）

\choices
    {\(-8\)}
    {\(-6\)}
    {\(6\)}
    {\(8\)}
\end{problem}

\begin{answer}
D
\end{answer}

\begin{problem}
圆 \(x^2+y^2-2x-8y+13=0\) 的圆心到直线 \(ax+y-1=0\) 的距离为 \(1\)，则 \(a=\)（\quad）

\choices
    {\(-\dfrac43\)}
    {\(-\dfrac34\)}
    {\(\sqrt3\)}
    {\(2\)}
\end{problem}

\begin{answer}
A
\end{answer}

\begin{problem}
如图，小明从街道的 \(E\) 处出发，先到 \(F\) 处与小红会合，再一起到位于 \(G\) 处的老年公寓参加志愿者活动，则小明到老年公寓可以选择的最短路径条数为（\quad）

\begingroup
\bitmapfigure[max width=0.70\linewidth]{img/2016/national_paper_2_science/q5_fig1.png}
\endgroup

\choices
    {\(24\)}
    {\(18\)}
    {\(12\)}
    {\(9\)}
\end{problem}

\begin{answer}
B
\end{answer}

\begin{problem}
如图是由圆柱与圆锥组合而成的几何体的三视图，则该几何体的表面积为（\quad）

\begingroup
\bitmapfigure[max width=0.42\linewidth]{img/2016/national_paper_2_science/q6_fig1.png}
\endgroup

\choices
    {\(20\pi\)}
    {\(24\pi\)}
    {\(28\pi\)}
    {\(32\pi\)}
\end{problem}

\begin{answer}
C
\end{answer}

\begin{problem}
若将函数 \(y=2\sin2x\) 的图象向左平移 \(\dfrac{\pi}{12}\) 个单位长度，则平移后的图象的对称轴为（\quad）

\choices
    {\(x=\dfrac{k\pi}{2}-\dfrac{\pi}{6}\)（\(k\in\Z\)）}
    {\(x=\dfrac{k\pi}{2}+\dfrac{\pi}{6}\)（\(k\in\Z\)）}
    {\(x=\dfrac{k\pi}{2}-\dfrac{\pi}{12}\)（\(k\in\Z\)）}
    {\(x=\dfrac{k\pi}{2}+\dfrac{\pi}{12}\)（\(k\in\Z\)）}
\end{problem}

\begin{answer}
B
\end{answer}

\begin{problem}
中国古代有计算多项式值的秦九韶算法，如图是实现该算法的程序框图. 执行该程序框图，若输入的 \(x=2\)，\(n=2\)，依次输入的 \(a\) 为 \(2\)，\(2\)，\(5\)，则输出的 \(s=\)（\quad）

\begingroup
\bitmapfigure[max width=6.0cm]{img/2016/national_paper_2_science/q8_fig1.png}
\endgroup

\choices
    {\(7\)}
    {\(12\)}
    {\(17\)}
    {\(34\)}
\end{problem}

\begin{answer}
C
\end{answer}

\begin{problem}
若 \(\cos\!\left(\dfrac{\pi}{4}-\alpha\right)=\dfrac35\)，则 \(\sin2\alpha=\)（\quad）

\choices
    {\(\dfrac{7}{25}\)}
    {\(\dfrac15\)}
    {\(-\dfrac15\)}
    {\(-\dfrac{7}{25}\)}
\end{problem}

\begin{answer}
D
\end{answer}

\begin{problem}
从区间 \([0,1]\) 随机抽取 \(2n\) 个数 \(x_1\)，\(x_2\)，\(\dots\)，\(x_n\)，\(y_1\)，\(y_2\)，\(\dots\)，\(y_n\).

将其构成 \(n\) 个数对 \((x_1,y_1)\)，\((x_2,y_2)\)，\(\dots\)，\((x_n,y_n)\).

其中两数的平方和小于 \(1\) 的数对共有 \(m\) 个，则用随机模拟的方法得到的圆周率 \(\pi\) 的近似值为（\quad）

\choices
    {\(\dfrac{4n}{m}\)}
    {\(\dfrac{2n}{m}\)}
    {\(\dfrac{4m}{n}\)}
    {\(\dfrac{2m}{n}\)}
\end{problem}

\begin{answer}
C
\end{answer}

\begin{problem}
已知 \(F_1\)，\(F_2\) 是双曲线 \(E:\dfrac{x^2}{a^2}-\dfrac{y^2}{b^2}=1\) 的左，右焦点，点 \(M\) 在 \(E\) 上，\(MF_1\perp x\) 轴，\(\sin\angle MF_2F_1=\dfrac13\)，则 \(E\) 的离心率为（\quad）

\choices
    {\(\sqrt2\)}
    {\(\dfrac32\)}
    {\(\sqrt3\)}
    {\(2\)}
\end{problem}

\begin{answer}
A
\end{answer}

\begin{problem}
已知函数 \(f(x)\)（\(x\in\R\)）满足 \(f(-x)=2-f(x)\)，若函数 \(y=\dfrac{x+1}{x}\) 与 \(y=f(x)\) 图象的交点为 \((x_1,y_1)\)，\((x_2,y_2)\)，\(\dots\)，\((x_m,y_m)\)，则 \(\displaystyle\sum_{i=1}^m(x_i+y_i)=\)（\quad）

\choices
    {\(0\)}
    {\(m\)}
    {\(2m\)}
    {\(4m\)}
\end{problem}

\begin{answer}
B
\end{answer}

\begin{solution}
记 \(g(x)=\dfrac{x+1}{x}=1+\dfrac1x\). 若 \((x_i,y_i)\) 是两个函数图象的交点，则 \(x_i\ne0\)，且 \(y_i=g(x_i)\). 由题设，\(f(-x_i)=2-f(x_i)=2-y_i\)，另一方面

\[
g(-x_i)=1-\dfrac1{x_i}=2-\left(1+\dfrac1{x_i}\right)=2-y_i.
\]

所以 \((-x_i,2-y_i)\) 也是两个函数图象的交点，并且它与 \((x_i,y_i)\) 不同. 这些交点可两两配对，每一对的横、纵坐标之和为

\[
x_i+(-x_i)+y_i+(2-y_i)=2.
\]

因此

\[
\sum_{i=1}^m(x_i+y_i)=\frac m2\cdot2=m.
\]

故选 B.
\end{solution}

\section{填空题}

\begin{problem}
\(\triangle ABC\) 的内角 \(A\)，\(B\)，\(C\) 的对边分别为 \(a\)，\(b\)，\(c\)，若 \(\cos A=\dfrac45\)，\(\cos C=\dfrac{5}{13}\)，\(a=1\)，则 \(b=\fillinblank{}\).
\end{problem}

\begin{answer}
\(\dfrac{21}{13}\)
\end{answer}


\begin{problem}
\(\alpha\)，\(\beta\) 是两个平面，\(m\)，\(n\) 是两条直线，有下列四个命题：
\begin{circlelist}
\item 如果 \(m\perp n\)，\(m\perp\alpha\)，\(n\parallel\beta\)，那么 \(\alpha\perp\beta\)；
\item 如果 \(m\perp\alpha\)，\(n\parallel\alpha\)，那么 \(m\perp n\)；
\item 如果 \(\alpha\parallel\beta\)，\(m\subset\alpha\)，那么 \(m\parallel\beta\)；
\item 如果 \(m\parallel n\)，\(\alpha\parallel\beta\)，那么 \(m\) 与 \(\alpha\) 所成的角和 \(n\) 与 \(\beta\) 所成的角相等.
\end{circlelist}
其中正确的命题是\fillinblank{}.
\end{problem}

\begin{answer}
\(\circled{2}\circled{3}\circled{4}\)
\end{answer}


\begin{problem}
有三张卡片，分别写有 \(1\) 和 \(2\)，\(1\) 和 \(3\)，\(2\) 和 \(3\). 甲，乙，丙三人各取走一张卡片，甲看了乙的卡片后说：“我与乙的卡片上相同的数字不是 \(2\)”；乙看了丙的卡片后说：“我与丙的卡片上相同的数字不是 \(1\)”；丙说：“我的卡片上的数字之和不是 \(5\)”；则甲的卡片上的数字是\fillinblank{}.
\end{problem}

\begin{answer}
\(1\) 和 \(3\)
\end{answer}


\begin{problem}
若直线 \(y=kx+b\) 是曲线 \(y=\ln x+2\) 的切线，也是曲线 \(y=\ln(x+1)\) 的切线，则 \(b=\fillinblank{}\).
\end{problem}

\begin{answer}
\(1-\ln2\)
\end{answer}

\begin{solution}
设直线与曲线 \(y=\ln x+2\) 的切点横坐标为 \(u\)，则 \(u>0\). 由导数，切线斜率为 \(1/u\)，切线方程为

\[
y=\frac1u x+\ln u+1.
\]

设直线与曲线 \(y=\ln(x+1)\) 的切点横坐标为 \(v\)，则 \(v>-1\). 该曲线在 \(v\) 处的切线方程为

\[
y=\frac1{v+1}x+\ln(v+1)-\frac v{v+1}.
\]

两条切线是同一直线，比较斜率得 \(u=v+1\). 再比较截距，得

\[
\ln u+1=\ln u-\frac{u-1}{u}.
\]

于是 \(2=1/u\)，所以 \(u=1/2\)，从而

\[
b=\ln u+1=1-\ln2.
\]
\end{solution}

\section{解答题}

\begin{problem}
\(S_n\) 为等差数列 \(\{a_n\}\) 的前 \(n\) 项和，且 \(a_1=1\)，\(S_7=28\)，记 \(b_n=[\lg a_n]\)，其中 \([x]\) 表示不超过 \(x\) 的最大整数，如 \([0.9]=0\)，\([\lg99]=1\).

\begin{enumerate}
\item 求 \(b_1\)，\(b_{11}\)，\(b_{101}\)；
\item 求数列 \(\{b_n\}\) 的前 \(1000\) 项和.
\end{enumerate}
\end{problem}

\begin{solution}
\begin{enumerate}
\item 设等差数列 \(\{a_n\}\) 的公差为 \(d\). 由 \(S_7=28\) 得

\[
28=\frac72\bigl(2a_1+6d\bigr)=7(1+3d).
\]

所以 \(d=1\)，从而 \(a_n=a_1+(n-1)d=n\). 因此

\[
b_1=\lfloor\lg1\rfloor=0,\qquad b_{11}=\lfloor\lg11\rfloor=1,\qquad b_{101}=\lfloor\lg101\rfloor=2.
\]

\item 当 \(1\leq n\leq9\) 时，\(0\leq\lg n<1\)，所以 \(b_n=0\)；当 \(10\leq n\leq99\) 时，\(1\leq\lg n<2\)，所以 \(b_n=1\)；当 \(100\leq n\leq999\) 时，\(2\leq\lg n<3\)，所以 \(b_n=2\)；最后 \(b_{1000}=3\). 因而

\[
\sum_{n=1}^{1000}b_n=9\times0+90\times1+900\times2+1\times3=1893.
\]
\end{enumerate}
\end{solution}


\begin{problem}
某险种的基本保费为 \(a\)（单位：元），继续购买该险种的投保人称为续保人，续保人本年度的保费与其上年度出险次数的关联如下：

\begin{center}
\begin{tabular}{c|cccccc}
上年度出险次数 & \(0\) & \(1\) & \(2\) & \(3\) & \(4\) & \(\geq 5\)\\
\hline
保费 & \(0.85a\) & \(a\) & \(1.25a\) & \(1.5a\) & \(1.75a\) & \(2a\)
\end{tabular}
\end{center}

设该险种一续保人一年内出险次数与相应概率如下：

\begin{center}
\begin{tabular}{c|cccccc}
一年内出险次数 & \(0\) & \(1\) & \(2\) & \(3\) & \(4\) & \(\geq 5\)\\
\hline
概率 & \(0.30\) & \(0.15\) & \(0.20\) & \(0.20\) & \(0.10\) & \(0.05\)
\end{tabular}
\end{center}

\begin{enumerate}
\item 求一续保人本年度的保费高于基本保费的概率；
\item 若一续保人本年度的保费高于基本保费，求其保费比基本保费高出 \(60\%\) 的概率；
\item 求续保人本年度的平均保费与基本保费的比值.
\end{enumerate}
\end{problem}

\begin{solution}
\begin{enumerate}
\item 保费高于基本保费对应上年度出险次数为 \(2\)，\(3\)，\(4\) 或不少于 \(5\) 次，所以所求概率为

\[
0.20+0.20+0.10+0.05=0.55.
\]

\item 保费比基本保费高出 \(60\%\) 或更多时，保费至少为 \(1.6a\)，对应出险次数为 \(4\) 或不少于 \(5\) 次. 因此在保费高于基本保费的条件下，所求条件概率为

\[
\frac{0.10+0.05}{0.55}=\frac{0.15}{0.55}=\frac3{11}.
\]

\item 以基本保费为单位，续保人本年度保费与基本保费的比值的期望为

\[
0.85\times0.30+1\times0.15+1.25\times0.20+1.5\times0.20+1.75\times0.10+2\times0.05=1.23.
\]
\end{enumerate}
\end{solution}


\begin{problem}
如图，菱形 \(ABCD\) 的对角线 \(AC\) 与 \(BD\) 交于点 \(O\)，\(AB=5\)，\(AC=6\)，点 \(E\)，\(F\) 分别在 \(AD\)，\(CD\) 上，\(AE=CF=\dfrac54\)，\(EF\) 交 \(BD\) 于点 \(H\)，将 \(\triangle DEF\) 沿 \(EF\) 折到 \(\triangle D'EF\) 的位置，\(OD'=\sqrt{10}\).
\begin{enumerate}
\item 证明：\(D'H\perp\) 平面 \(ABCD\)；
\item 求二面角 \(B-D'A-C\) 的正弦值.
\end{enumerate}

\begingroup
\bitmapfigure[max width=0.34\linewidth]{img/2016/national_paper_2_science/q19_fig1.png}
\endgroup
\end{problem}

\begin{solution}
\begin{enumerate}
\item 在平面 \(ABCD\) 内建立直角坐标系，取

\[
A=(-3,0),\qquad B=(0,-4),\qquad C=(3,0),\qquad D=(0,4).
\]

因为 \(AD=CD=AB=5\)，且 \(AE=CF=5/4\)，所以

\[
E=\left(-\frac94,1\right),\qquad F=\left(\frac94,1\right).
\]

于是 \(EF\) 的方程为 \(y=1\)，\(BD\) 的方程为 \(x=0\)，故

\[
H=(0,1),\qquad OH=1,\qquad DH=3.
\]

又 \(DE=DF\)，折叠后 \(D'E=D'F\)，且 \(H\) 是 \(EF\) 的中点，所以在等腰三角形 \(D'EF\) 中，\(D'H\perp EF\). 折叠不改变点 \(H\) 到顶点的距离，故 \(HD'=HD=3\). 结合 \(OD'=\sqrt{10}\)，有

\[
OD'^2=10=1+9=OH^2+HD'^2.
\]

由勾股定理的逆定理，\(D'H\perp OH\)，从而 \(D'H\perp BD\). 因为 \(D'H\) 同时垂直于平面 \(ABCD\) 内相交的两条直线 \(BD\)，\(EF\)，所以 \(D'H\perp\) 平面 \(ABCD\).

\item 由上面的坐标和 \(HD'=3\)，可取空间直角坐标系使

\[
D'=(0,1,3).
\]

于是

\[
\overrightarrow{AB}=(3,-4,0),\qquad \overrightarrow{AC}=(6,0,0),\qquad \overrightarrow{AD'}=(3,1,3).
\]

取平面 \(D'AB\)，\(D'AC\) 的法向量分别为

\[
\bs n_1=\overrightarrow{AB}\times\overrightarrow{AD'}=(-12,-9,15),
\]

\[
\bs n_2=\overrightarrow{AC}\times\overrightarrow{AD'}=(0,-18,6).
\]

因此两平面所成锐二面角的余弦为

\[
\cos\theta=\frac{|\bs n_1\cdot\bs n_2|}{|\bs n_1|\,|\bs n_2|}=\frac{252}{15\sqrt2\cdot6\sqrt{10}}=\frac7{5\sqrt5}.
\]

所以

\[
\sin\theta=\sqrt{1-\frac{49}{125}}=\frac{2\sqrt{95}}{25}.
\]
\end{enumerate}
\end{solution}


\begin{problem}
已知椭圆 \(E:\dfrac{x^2}{t}+\dfrac{y^2}{3}=1\)，焦点在 \(x\) 轴上，\(A\) 是 \(E\) 的左顶点，斜率为 \(k\)（\(k>0\)）的直线交 \(E\) 于 \(A\)，\(M\) 两点，点 \(N\) 在 \(E\) 上，且 \(MA\perp NA\).
\begin{enumerate}
\item 当 \(t=4\)，\(\lvert AM\rvert=\lvert AN\rvert\) 时，求 \(\triangle AMN\) 的面积；
\item 当 \(2\lvert AM\rvert=\lvert AN\rvert\) 时，求 \(k\) 的取值范围.
\end{enumerate}
\end{problem}

\begin{solution}
\begin{enumerate}
\item 记椭圆的半长轴、半短轴分别为 \(a_0=\sqrt t\)，\(b_0=\sqrt3\)，则左顶点 \(A=(-a_0,0)\). 对于过 \(A\) 且斜率为 \(s\) 的直线，有 \(y=s(x+a_0)\). 代入椭圆方程，除去交点 \(A\) 对应的根后，得到另一个交点 \(P_s\) 满足

\[
x_{P_s}+a_0=\frac{2a_0b_0^2}{b_0^2+a_0^2s^2}.
\]

所以

\[
AP_s=\frac{2a_0b_0^2\sqrt{1+s^2}}{b_0^2+a_0^2s^2}.
\]

取 \(s=k\) 得点 \(M\)，取 \(s=-1/k\) 得点 \(N\)，从而

\[
\frac{|AN|}{|AM|}=\frac{k\left(b_0^2+a_0^2k^2\right)}{a_0^2+b_0^2k^2}.
\]

当 \(t=4\) 时，\(a_0^2=4\)，\(b_0^2=3\). 由 \(|AM|=|AN|\) 得

\[
k(3+4k^2)=4+3k^2.
\]

即

\[
(k-1)(4k^2+k+4)=0.
\]

因为 \(4k^2+k+4>0\)，且 \(k>0\)，所以 \(k=1\). 此时

\[
|AM|=\frac{2\cdot2\cdot3\sqrt2}{3+4}=\frac{12\sqrt2}{7}.
\]

又 \(MA\perp NA\)，所以

\[
S_{\triangle AMN}=\frac12|AM|\,|AN|=\frac12\left(\frac{12\sqrt2}{7}\right)^2=\frac{144}{49}.
\]

\item 由上面的长度比公式和 \(a_0^2=t\)，\(b_0^2=3\)，条件 \(2|AM|=|AN|\) 等价于

\[
\frac{k(3+tk^2)}{t+3k^2}=2.
\]

整理得

\[
t(k^3-2)=3k(2k-1).
\]

由于椭圆的焦点在 \(x\) 轴上，必须有 \(t>3\). 由上式

\[
t-3=\frac{3(2-k)(k^2+1)}{k^3-2}.
\]

因此 \(t>3\) 等价于

\[
\frac{2-k}{k^3-2}>0.
\]

结合 \(k>0\) 可得 \(k^3>2\) 且 \(k<2\)，即

\[
k\in\left(\sqrt[3]{2},2\right).
\]
\end{enumerate}
\end{solution}


\begin{problem}
\begin{enumerate}
\item 讨论函数 \(f(x)=\dfrac{x-2}{x+2}\e^x\) 的单调性，并证明当 \(x>0\) 时，\((x-2)\e^x+x+2>0\).
\item 证明：当 \(a\in[0,1)\) 时，函数 \(g(x)=\dfrac{\e^x-ax-a}{x^2}\)（\(x>0\)）有最小值. 设 \(g(x)\) 的最小值为 \(h(a)\)，求函数 \(h(a)\) 的值域.
\end{enumerate}
\end{problem}

\begin{solution}
\begin{enumerate}
\item 函数的定义域为 \(\R\) 中除 \(x=-2\) 外的所有数. 求导得

\[
f'(x)=\e^x\left(\frac{x-2}{x+2}+\frac4{(x+2)^2}\right)=\frac{x^2\e^x}{(x+2)^2}.
\]

当 \(x\ne-2\) 时，\(f'(x)\geq0\)，且除 \(x=0\) 外为正，所以 \(f(x)\) 在 \((-\infty,-2)\) 和 \((-2,+\infty)\) 上均单调递增.

令 \(p(x)=(x-2)\e^x+x+2\). 则 \(p(0)=0\)，且

\[
p'(x)=(x-1)\e^x+1,\qquad p''(x)=x\e^x.
\]

当 \(x>0\) 时，\(p''(x)>0\)，所以 \(p'(x)>p'(0)=0\)，进而 \(p(x)>p(0)=0\). 即当 \(x>0\) 时，\((x-2)\e^x+x+2>0\).

\item 记

\[
q_a(x)=(x-2)\e^x+a(x+2).
\]

则

\[
g'(x)=\frac{q_a(x)}{x^3}.
\]

在 \(0\leq x\leq2\) 上令

\[
\varphi(x)=\frac{(2-x)\e^x}{x+2}.
\]

当 \(0<x<2\) 时，

\[
\varphi'(x)=-\frac{x^2\e^x}{(x+2)^2}<0.
\]

又 \(\varphi(0)=1\)，\(\varphi(2)=0\)，所以对任意 \(a\in[0,1)\)，存在唯一的 \(x_a\in(0,2]\)，使得 \(a=\varphi(x_a)\). 对 \(0<x<2\)，有

\[
q_a(x)=(x+2)\bigl(a-\varphi(x)\bigr).
\]

由 \(\varphi\) 的单调性，\(q_a(x)<0\)（\(0<x<x_a\)），\(q_a(x)>0\)（\(x_a<x<2\)）；当 \(x>2\) 时，\(q_a(x)>0\). 因此 \(g(x)\) 先减后增，在 \(x=x_a\) 处取得最小值.

由 \(q_a(x_a)=0\) 得 \(a=(2-x_a)\e^{x_a}/(x_a+2)\)，代入 \(g\) 可得

\[
h(a)=g(x_a)=\frac{\e^{x_a}-a(x_a+1)}{x_a^2}=\frac{\e^{x_a}}{x_a+2}.
\]

当 \(x>0\) 时，函数 \(r(x)=\dfrac{\e^x}{x+2}\) 满足

\[
r'(x)=\frac{\e^x(x+1)}{(x+2)^2}>0.
\]

而 \(a\in[0,1)\) 对应 \(x_a\in(0,2]\)，所以

\[
h(a)\in\left(\frac12,\frac{\e^2}{4}\right].
\]
\end{enumerate}
\end{solution}


\section{选考题：解答题}

\begin{problem}
如图，在正方形 \(ABCD\) 中，\(E\)，\(G\) 分别在边 \(DA\)，\(DC\) 上（不与端点重合），且 \(DE=DG\)，过 \(D\) 点作 \(DF\perp CE\)，垂足为 \(F\).
\begin{enumerate}
\item 证明：\(B\)，\(C\)，\(G\)，\(F\) 四点共圆；
\item 若 \(AB=1\)，\(E\) 为 \(DA\) 的中点，求四边形 \(BCGF\) 的面积.
\end{enumerate}

\begingroup
\bitmapfigure[max width=0.34\linewidth]{img/2016/national_paper_2_science/q22_fig1.png}
\endgroup
\end{problem}

\begin{solution}
\begin{enumerate}
\item 设正方形边长为 \(s\)，令 \(r=DE/s=DG/s\)，其中 \(0<r<1\). 建立直角坐标系，取

\[
D=(0,0),\qquad C=(s,0),\qquad A=(0,-s),\qquad B=(s,-s).
\]

则 \(E=(0,-rs)\)，\(G=(rs,0)\). 直线 \(CE\) 的斜率为 \(r\)，所以过 \(D\) 且垂直于 \(CE\) 的直线为 \(y=-x/r\). 解得垂足

\[
F=\left(\frac{sr^2}{1+r^2},-\frac{sr}{1+r^2}\right).
\]

由此

\[
\overrightarrow{FB}=\frac{s}{1+r^2}\bigl(1,-1-r^2+r\bigr),
\]

\[
\overrightarrow{FG}=\frac{s}{1+r^2}\bigl(r(1-r+r^2),r\bigr).
\]

于是

\[
\overrightarrow{FB}\cdot\overrightarrow{FG}
=\frac{s^2}{(1+r^2)^2}\left(r(1-r+r^2)+r(-1-r^2+r)\right)=0.
\]

即 \(BF\perp GF\). 又 \(BC\perp CG\)，所以

\[
\angle BFG=\angle BCG=90^\circ.
\]

两角同对弦 \(BG\)，故 \(B\)，\(C\)，\(G\)，\(F\) 四点共圆.

\item 当 \(AB=1\) 且 \(E\) 为 \(DA\) 的中点时，取

\[
B=(1,-1),\qquad C=(1,0),\qquad G=\left(\frac12,0\right),\qquad F=\left(\frac15,-\frac25\right).
\]

此时 \(BC=BF=1\)，\(CG=GF=1/2\)，且 \(\triangle BCG\) 与 \(\triangle BFG\) 都是直角三角形，斜边同为 \(BG\)，所以两三角形全等. 因而

\[
S_{BCGF}=2S_{\triangle BCG}=2\cdot\frac12\cdot1\cdot\frac12=\frac12.
\]
\end{enumerate}
\end{solution}


\begin{problem}
在直角坐标系 \(xOy\) 中，圆 \(C\) 的方程为 \((x+6)^2+y^2=25\).
\begin{enumerate}
\item 以坐标原点为极点，\(x\) 轴正半轴为极轴建立极坐标系，求 \(C\) 的极坐标方程；
\item 直线 \(l\) 的参数方程是
\(
\begin{cases}
x=t\cos\alpha,\\
y=t\sin\alpha,
\end{cases}
\)
（\(t\) 为参数），\(l\) 与 \(C\) 交于 \(A\)，\(B\) 两点，\(\lvert AB\rvert=\sqrt{10}\)，求 \(l\) 的斜率.
\end{enumerate}
\end{problem}

\begin{solution}
\begin{enumerate}
\item 将 \(x=\rho\cos\theta\)，\(y=\rho\sin\theta\) 代入圆方程，得

\[
\rho^2\cos^2\theta+12\rho\cos\theta+11+\rho^2\sin^2\theta=0.
\]

利用 \(\sin^2\theta+\cos^2\theta=1\)，所以圆 \(C\) 的极坐标方程为

\[
\rho^2+12\rho\cos\theta+11=0.
\]

\item 圆心为 \(P=(-6,0)\)，半径为 \(5\). 直线 \(l\) 的方向向量为 \((\cos\alpha,\sin\alpha)\)，其方程可写为

\[
x\sin\alpha-y\cos\alpha=0.
\]

所以圆心到直线 \(l\) 的距离为 \(6|\sin\alpha|\). 由弦长公式

\[
|AB|=2\sqrt{25-36\sin^2\alpha}=\sqrt{10}.
\]

两边平方得 \(100-144\sin^2\alpha=10\)，即 \(\sin^2\alpha=5/8\). 因而 \(\cos^2\alpha=3/8\)，从而直线斜率满足

\[
k^2=\tan^2\alpha=\frac{\sin^2\alpha}{\cos^2\alpha}=\frac53.
\]

故

\[
k=\pm\frac{\sqrt{15}}3.
\]
\end{enumerate}
\end{solution}


\begin{problem}
已知函数 \(f(x)=\lvert x-\dfrac12\rvert+\lvert x+\dfrac12\rvert\)，\(M\) 为不等式 \(f(x)<2\) 的解集.
\begin{enumerate}
\item 求 \(M\)；
\item 证明：当 \(a\)，\(b\in M\) 时，\(\lvert a+b\rvert<\lvert 1+ab\rvert\).
\end{enumerate}
\end{problem}

\begin{solution}
\begin{enumerate}
\item 当 \(x\leq-1/2\) 时，\(f(x)=-2x\)，不等式 \(f(x)<2\) 等价于 \(-1<x\leq-1/2\)；当 \(-1/2\leq x\leq1/2\) 时，\(f(x)=1<2\)；当 \(x\geq1/2\) 时，\(f(x)=2x\)，不等式等价于 \(1/2\leq x<1\). 因此

\[
M=(-1,1).
\]

\item 因为 \(a,b\in(-1,1)\)，所以 \(1-a^2>0\)，\(1-b^2>0\)，并且 \(ab>-1\)，从而 \(1+ab>0\). 又

\[
\left(1+ab\right)^2-(a+b)^2=(1-a^2)(1-b^2)>0.
\]

因此 \(1+ab>|a+b|\)，即

\[
|a+b|<|1+ab|.
\]
\end{enumerate}
\end{solution}
